% Generated draft: requires mathematical and source review.
\documentclass[11pt]{article}
\usepackage[margin=25mm]{geometry}
\usepackage{fontspec}
\setmainfont{Latin Modern Roman}
\usepackage{amsmath,amssymb,mathtools,unicode-math}
\setmathfont{Latin Modern Math}
\usepackage{enumitem,needspace,xurl,hyperref}
\setlength{\parindent}{0pt}
\setlength{\parskip}{5pt}
\setlength{\emergencystretch}{3em}
\newcommand{\N}{\mathbb{N}}
\newcommand{\Z}{\mathbb{Z}}
\newcommand{\Q}{\mathbb{Q}}
\newcommand{\R}{\mathbb{R}}
\newcommand{\C}{\mathbb{C}}
\newcommand{\F}{\mathbb{F}}
\newcommand{\A}{\mathbb{A}}
\newcommand{\PP}{\mathbb{P}}
\newcommand{\E}{\mathbb{E}}
\newcommand{\eps}{\varepsilon}
\newcommand{\dd}{\,\mathrm{d}}
\begin{document}
\section*{UnsolvedMath ID: 30001166}
\textbf{Source problem code:} catalogue entry
\subsection{Definitions and mathematical statement}
The supplied record concerns magnetic Lieb--Thirring inequalities, but does not specify the operator class, regularity assumptions, or the precise comparison constant. The following is an explicit conventional scalar formulation on smooth Euclidean test data; extending it to the full intended class requires clarification.

Fix an integer \(n\geq 1\) and \(1\leq\gamma<3/2\). Let \(A\in C^\infty(\R^n;\R^n)\) be a real magnetic vector potential, with magnetic field represented by \(B_{jk}=\partial_jA_k-\partial_kA_j\). For a real potential \(V\in C_c^\infty(\R^n)\), let
\[
H_{A,V}=(-i\nabla-A)^2+V
\]
be the self-adjoint operator obtained from the Friedrichs extension of its quadratic form initially defined on \(C_c^\infty(\R^n)\). Write \(V_-(x)=\max\{-V(x),0\}\) and
\[
S_\gamma(A,V)=\operatorname{Tr}(H_{A,V})_-^\gamma.
\]
Here the negative part is defined by spectral calculus; when the negative spectrum is discrete, this trace is the sum of \(|\lambda|^\gamma\) over negative eigenvalues, counted with multiplicity. The trace is allowed to be infinite.

For this test class define
\[
L_{\gamma,n}^{\mathrm{mag}}
=\sup_{\substack{A\in C^\infty(\R^n;\R^n),\ V\in C_c^\infty(\R^n;\R)\\ \int_{\R^n}V_-^{\gamma+n/2}\,\dd x>0}}
\frac{S_\gamma(A,V)}{\int_{\R^n}V_-^{\gamma+n/2}\,\dd x},
\qquad
L_{\gamma,n}^{0}
=\sup_{\substack{V\in C_c^\infty(\R^n;\R)\\ \int_{\R^n}V_-^{\gamma+n/2}\,\dd x>0}}
\frac{S_\gamma(0,V)}{\int_{\R^n}V_-^{\gamma+n/2}\,\dd x}.
\]
The Euclidean questions are whether a finite constant independent of \(A\) and \(V\) gives
\[
S_\gamma(A,V)\leq L_{\gamma,n}\int_{\R^n}V_-^{\gamma+n/2}\,\dd x,
\]
and, more specifically, whether magnetic fields alter the optimal constant, expressed in this formulation as the question
\[
L_{\gamma,n}^{\mathrm{mag}}\stackrel{?}{=}L_{\gamma,n}^{0}.
\]
The original statement particularly asks whether, for \(n=3\) and \(\gamma=1\), there are magnetic and electric potentials violating the intended classical comparison bound.

The phrase ``classical optimal constant'' is unresolved in the supplied record. If it means the optimal nonmagnetic constant, the proposed counterexample would satisfy
\[
S_1(A,V)>L_{1,3}^{0}\int_{\R^3}V_-^{5/2}\,\dd x.
\]
If it instead means the semiclassical phase-space constant
\[
L_{\gamma,n}^{\mathrm{cl}}=(2\pi)^{-n}\int_{\R^n}(1-|p|^2)_+^\gamma\,\dd p,
\qquad t_+=\max\{t,0\},
\]
then the proposed counterexample would use \(L_{1,3}^{\mathrm{cl}}\) in that inequality. These two comparison targets must not be identified without further source clarification. The record also does not specify whether ``Li--Yau/Lieb--Thirring bound'' refers to this negative-eigenvalue moment inequality, an eigenvalue-sum formulation, or a particular sharp constant.

The geometric question asks for Lieb--Thirring inequalities on Riemannian manifolds and for the dependence of their critical exponents on geometry. A conventional candidate, after fixing an \(n\)-dimensional Riemannian manifold \((M,g)\), a self-adjoint realization, and admissible potentials, is
\[
\operatorname{Tr}(-\Delta_g+V)_-^\gamma
\leq C_{\gamma,M,g}\int_M V_-^{\gamma+n/2}\,\dd\operatorname{vol}_g,
\]
with \(C_{\gamma,M,g}<\infty\) independent of \(V\). A magnetic version would replace \(-\Delta_g\) by a specified magnetic Laplacian and would require an explicit decision about uniformity in the magnetic field. One can then ask which \(\gamma\) admit the chosen inequality, and how that set depends on geometry. The supplied record does not determine the manifold class, boundary conditions, admissible operator and potential classes, allowable geometric dependence of constants, or whether additional terms are permitted. Consequently it does not determine a unique manifold inequality or a unique meaning of its critical exponent.

The supplied dated assessment reports universal magnetic comparison bounds and identifies sharp magnetic constants, including the \(n=3,\gamma=1\) counterexample question, as unresolved. This is a supplied qualification, not an independently verified mathematical or literature conclusion.
\subsection{Short English statement}
The problem asks whether arbitrary magnetic fields preserve the intended Lieb--Thirring bound and its optimal constant for \(1\leq\gamma<3/2\), especially when \(n=3\) and \(\gamma=1\). It also asks which analogous inequalities hold on manifolds and how their critical exponents depend on geometry. The comparison constant and several analytic and geometric hypotheses remain unspecified; the displayed smooth scalar formulation is conditional and requires human review.
\subsection{Sources}
\begin{itemize}
\item Supplied local UnsolvedMath registry record: permanent ID \(30001166\), exact source problem code \texttt{catalogue entry}, title ``Magnetic Lieb--Thirring Constants and Geometric Critical Exponents,'' in the Oberwolfach Reports Open Problems collection. The proposed year is supplied as 2009; no individual proposer is named.
\item Supplied report citation: \emph{Low Eigenvalues of Laplace and Schr\"odinger Operators} (2009), Oberwolfach Reports, pp.\ 355--428, DOI \texttt{10.4171/owr/2009/06}. \url{https://doi.org/10.4171/owr/2009/06}. The background describes this report as documenting workshop discussions and open problems. No problem-specific page or speaker attribution is supplied.
\item Supplied literature source, labelled ``magnetic Lieb--Thirring inequalities'': \url{https://doi.org/10.1007/s13373-015-0067-9}.
\item Supplied literature source, labelled ``best Lieb--Thirring constants numerical investigation'': \url{https://arxiv.org/abs/1206.1473}.
\item Literature assessment date: \texttt{2026-08-21}. This is supplied dated curation, not independently verified. It labels the problem open, reports universal magnetic comparison bounds, and describes the sharp magnetic constant in the stated range, including the \(n=3,\gamma=1\) counterexample question, as unresolved. The background expressly requests independent verification of time-sensitive status and cited claims.
\item The registry supplies statement status \texttt{corrected\_verified} and says that the mathematical content was checked against the official report, with notation, context, and question form normalized. That verification claim is part of the supplied provenance; this draft does not independently confirm it.
\end{itemize}
\end{document}
